\begin{examplebox}
\textbf{例1：没有两条实切线时，只能称为极线。}
双曲线 $x^2/4-y^2=1$，$P=(3,1)$。其 $q=9/4-1=5/4>1$，
所以没有实切点弦。形式极线为 $3x-4y=4$；代入 $y=(3x-4)/4$，
得到 $5x^2-24x+32=0$，判别式 $576-640=-64<0$。
\medskip
\textbf{例2：纵轴双曲线也须检验。}
对 $y^2/9-x^2/4=1$，点 $P=(2,6)$ 的 $q=4-1=3>1$，
故没有两条实切线，也无实切点弦；其极线 $4y-3x=6$ 不能称作切点弦。
\medskip
\textbf{例3：一个实际存在的切点弦。}
对 $x^2/4-y^2=1$，取 $P=(0,2)$，$q=-4<0$，满足存在条件。
极线为 $-2y=1$，即 $y=-1/2$。与曲线联立得切点
$A=(\sqrt5,-1/2)$、$B=(-\sqrt5,-1/2)$。
两点处切线为 $\pm\sqrt5 x/4+y/2=1$，都经过 $P=(0,2)$。
\end{examplebox}
